% GENERATED FILE. Edit canonical lessons, then run npm run generate:books.
% canonical-source-hash: fnv1a64:62fc6dba79025c3c
% canonical-lessons: UR-C521-maich, UR-C521-turnament, UR-C521-kaptan, UR-C521-fatih, UR-C521-shikast

\chapter{The Big Match}
\label{ch:ur-the-big-match}
\hlchaptermodality{521}

\begin{chapteropening}
\textbf{By the end of this chapter:} \emph{I can talk about matches, winners and losers.}

Complete the last lesson of chapter 521: I can talk about matches, winners and losers.
\end{chapteropening}

\section[\texorpdfstring{\ur{میچ} (maich)}{میچ (maich)}]{\ur{میچ} (maich) — a match (a game)}

\label{lesson:UR-C521-maich}



\begin{quote}
Before the new one: say the Urdu for to pit against each other, then the Urdu for a loincloth.
\end{quote}

\subsection*{You'll want to know: \ur{میچ}}
\textbf{\ur{میچ}} — \emph{maich} — \textquotedblleft{}a match (a game)\textquotedblright{}.

\begin{quote}
\ur{کل} \ur{ہماری} \ur{ٹیم} \ur{کا} \ur{میچ} \ur{ہے۔}
\end{quote}

\begin{quote}
\emph{kal hamārī ṭīm kā maich hai.}
\end{quote}

\begin{quote}
Our team has a match tomorrow.
\end{quote}

Masculine.

\begin{cousinweb}
\textbf{\ur{میچ}} (\emph{maich}) is borrowed from English \emph{match}.
\end{cousinweb}

\begin{grammarlens}[title={What you've built}]
One more word for describing what you lived and felt.
\end{grammarlens}

\subsection*{Your turn}
\begin{itemize}
\raggedright
  \item \emph{Say it:} \emph{maich}
  \item \emph{Say it:} \emph{maich}, once more
  \item \emph{Say it:} \emph{kal hamārī ṭīm kā maich hai.} — Our team has a match tomorrow.
  \item \emph{Recall:} say the Urdu for to pit against each other, then the Urdu for a loincloth, then say \emph{maich} again
\end{itemize}

\begin{tcolorbox}[breakable,colback=teal!4,colframe=teal!35!black,title={Before you move on}]
What does \ur{میچ} mean? (\textquotedblleft{}A match (a game)\textquotedblright{}.) Say it once more.
\end{tcolorbox}
\section[\texorpdfstring{\ur{ٹورنامنٹ} (ṭūrnāmenṭ)}{ٹورنامنٹ (ṭūrnāmenṭ)}]{\ur{ٹورنامنٹ} (ṭūrnāmenṭ) — a tournament}

\label{lesson:UR-C521-turnament}



\begin{quote}
Before the new one: say the Urdu for a loincloth, then the Urdu for a match (a game).
\end{quote}

\subsection*{You'll want to know: \ur{ٹورنامنٹ}}
\textbf{\ur{ٹورنامنٹ}} — \emph{ṭūrnāmenṭ} — \textquotedblleft{}a tournament\textquotedblright{}.

\begin{quote}
\ur{ہماری} \ur{ٹیم} \ur{نے} \ur{ٹورنامنٹ} \ur{جیتا۔}
\end{quote}

\begin{quote}
\emph{hamārī ṭīm ne ṭūrnāmenṭ jītā.}
\end{quote}

\begin{quote}
Our team won the tournament.
\end{quote}

Masculine.

\begin{cousinweb}
\textbf{\ur{ٹورنامنٹ}} (\emph{ṭūrnāmenṭ}) is borrowed from English \emph{tournament}.
\end{cousinweb}

\begin{grammarlens}[title={What you've built}]
One more word for describing what you lived and felt.
\end{grammarlens}

\subsection*{Your turn}
\begin{itemize}
\raggedright
  \item \emph{Say it:} \emph{ṭūrnāmenṭ}
  \item \emph{Say it:} \emph{ṭūrnāmenṭ}, once more
  \item \emph{Say it:} \emph{hamārī ṭīm ne ṭūrnāmenṭ jītā.} — Our team won the tournament.
  \item \emph{Recall:} say the Urdu for a loincloth, then the Urdu for a match (a game), then say \emph{ṭūrnāmenṭ} again
\end{itemize}

\begin{tcolorbox}[breakable,colback=teal!4,colframe=teal!35!black,title={Before you move on}]
What does \ur{ٹورنامنٹ} mean? (\textquotedblleft{}A tournament\textquotedblright{}.) Say it once more.
\end{tcolorbox}
\section[\texorpdfstring{\ur{کپتان} (kaptān)}{کپتان (kaptān)}]{\ur{کپتان} (kaptān) — a captain (of a team)}

\label{lesson:UR-C521-kaptan}



\begin{quote}
Before the new one: say the Urdu for a match (a game), then the Urdu for a tournament.
\end{quote}

\subsection*{You'll want to know: \ur{کپتان}}
\textbf{\ur{کپتان}} — \emph{kaptān} — \textquotedblleft{}a captain (of a team)\textquotedblright{}.

\begin{quote}
\ur{کپتان} \ur{نے} \ur{کھلاڑیوں} \ur{کو} \ur{بلایا۔}
\end{quote}

\begin{quote}
\emph{kaptān ne khilāṛiyoṅ ko bulāyā.}
\end{quote}

\begin{quote}
The captain called the players over.
\end{quote}

Masculine or feminine, by the person meant.

\begin{cousinweb}
\textbf{\ur{کپتان}} (\emph{kaptān}) is usually traced to Portuguese \emph{capitão}, \textquotedblleft{}a captain\textquotedblright{}.
\end{cousinweb}

\begin{grammarlens}[title={What you've built}]
One more word for describing what you lived and felt.
\end{grammarlens}

\subsection*{Your turn}
\begin{itemize}
\raggedright
  \item \emph{Say it:} \emph{kaptān}
  \item \emph{Say it:} \emph{kaptān}, once more
  \item \emph{Say it:} \emph{kaptān ne khilāṛiyoṅ ko bulāyā.} — The captain called the players over.
  \item \emph{Recall:} say the Urdu for a match (a game), then the Urdu for a tournament, then say \emph{kaptān} again
\end{itemize}

\begin{tcolorbox}[breakable,colback=teal!4,colframe=teal!35!black,title={Before you move on}]
What does \ur{کپتان} mean? (\textquotedblleft{}A captain (of a team)\textquotedblright{}.) Say it once more.
\end{tcolorbox}
\section[\texorpdfstring{\ur{فاتح} (fātih)}{فاتح (fātih)}]{\ur{فاتح} (fātih) — a winner; a conqueror}

\label{lesson:UR-C521-fatih}



\begin{quote}
Before the new one: say the Urdu for a tournament, then the Urdu for a captain (of a team).
\end{quote}

\subsection*{You'll want to know: \ur{فاتح}}
\textbf{\ur{فاتح}} — \emph{fātih} — \textquotedblleft{}a winner; a conqueror\textquotedblright{}.

\begin{quote}
\ur{آخر} \ur{میں} \ur{ہماری} \ur{ٹیم} \ur{فاتح} \ur{رہی۔}
\end{quote}

\begin{quote}
\emph{ākhir meṅ hamārī ṭīm fātih rahī.}
\end{quote}

\begin{quote}
In the end our team came out the winner.
\end{quote}

Masculine or feminine, by the person meant.

\begin{cousinweb}
\textbf{\ur{فاتح}} (\emph{fātih}) is usually traced, through Persian, to Arabic \emph{fātiḥ}, \textquotedblleft{}one who opens; a conqueror\textquotedblright{}.
\end{cousinweb}

\begin{grammarlens}[title={What you've built}]
One more word for describing what you lived and felt.
\end{grammarlens}

\subsection*{Your turn}
\begin{itemize}
\raggedright
  \item \emph{Say it:} \emph{fātih}
  \item \emph{Say it:} \emph{fātih}, once more
  \item \emph{Say it:} \emph{ākhir meṅ hamārī ṭīm fātih rahī.} — In the end our team came out the winner.
  \item \emph{Recall:} say the Urdu for a tournament, then the Urdu for a captain (of a team), then say \emph{fātih} again
\end{itemize}

\begin{tcolorbox}[breakable,colback=teal!4,colframe=teal!35!black,title={Before you move on}]
What does \ur{فاتح} mean? (\textquotedblleft{}A winner; a conqueror\textquotedblright{}.) Say it once more.
\end{tcolorbox}
\section[\texorpdfstring{\ur{شکست} (shikast)}{شکست (shikast)}]{\ur{شکست} (shikast) — a defeat}

\label{lesson:UR-C521-shikast}



\begin{quote}
Before the new one: say the Urdu for a captain (of a team), then the Urdu for a winner.
\end{quote}

\subsection*{You'll want to know: \ur{شکست}}
\textbf{\ur{شکست}} — \emph{shikast} — \textquotedblleft{}a defeat\textquotedblright{}.

\begin{quote}
\ur{ہماری} \ur{ٹیم} \ur{کو} \ur{پہلے} \ur{میچ} \ur{میں} \ur{شکست} \ur{ہوئی۔}
\end{quote}

\begin{quote}
\emph{hamārī ṭīm ko pehle maich meṅ shikast huī.}
\end{quote}

\begin{quote}
Our team was beaten in the first match.
\end{quote}

Feminine.

\begin{cousinweb}
\textbf{\ur{شکست}} (\emph{shikast}) is usually traced to Persian \emph{shikast}, \textquotedblleft{}a breaking; a defeat\textquotedblright{}, from \emph{shikastan}, \textquotedblleft{}to break\textquotedblright{}.
\end{cousinweb}

\begin{grammarlens}[title={What you've built}]
One more word for describing what you lived and felt.
\end{grammarlens}

\subsection*{Your turn}
\begin{itemize}
\raggedright
  \item \emph{Say it:} \emph{shikast}
  \item \emph{Say it:} \emph{shikast}, once more
  \item \emph{Say it:} \emph{hamārī ṭīm ko pehle maich meṅ shikast huī.} — Our team was beaten in the first match.
  \item \emph{Recall:} say the Urdu for a captain (of a team), then the Urdu for a winner, then say \emph{shikast} again
\end{itemize}

\begin{tcolorbox}[breakable,colback=teal!4,colframe=teal!35!black,title={Before you move on}]
What does \ur{شکست} mean? (\textquotedblleft{}A defeat\textquotedblright{}.) Say it once more.
\end{tcolorbox}
